\section{Линейные непрерывные функционалы}
\label{sec:continuous-linear-functionals}

\subsection{Линейные функционалы}

Пусть $X$ --- линейное нормированное пространство над $\mathbb R$
(комплексный случай формулируется аналогично). \emph{Линейным функционалом} на $X$
называется линейное отображение
\[
    f\colon X\to\mathbb R.
\]
то есть
\[
    f(\alpha x+\beta y)=\alpha f(x)+\beta f(y)
\]
для любых $x,y\in X$ и $\alpha,\beta\in\mathbb R$.

Простейший пример в $\mathbb R^n$ --- скалярное произведение с фиксированным
вектором $a$:
\[
    f(x)=(x,a)=\sum_{k=1}^n a_kx_k.
\]

\subsection{Непрерывность и ограниченность}

Для линейного функционала непрерывность эквивалентна ограниченности. Функционал
$f$ называется \emph{ограниченным}, если существует константа $C\geq0$ такая,
что
\[
    |f(x)|\leq C\|x\|,
    \qquad \forall x\in X.
\]

Для линейного функционала эквивалентны следующие утверждения:
\begin{enumerate}
    \item $f$ непрерывен на $X$;
    \item $f$ непрерывен в нуле;
    \item $f$ ограничен.
\end{enumerate}
Действительно, из оценки ограниченности сразу следует
\[
    |f(x)-f(y)|=|f(x-y)|\leq C\|x-y\|,
\]
то есть непрерывность. Обратное утверждение следует из линейности и
непрерывности в нуле.

\subsection{Норма функционала}

Норма линейного непрерывного функционала определяется как
\[
    \|f\|_{X^*}
    =\sup_{\|x\|\leq1}|f(x)|.
\]
Эквивалентно, это наименьшая константа $C$, для которой
\[
    |f(x)|\leq C\|x\|\qquad\forall x\in X.
\]
Поэтому всегда справедлива фундаментальная оценка
\[
    |f(x)|\leq\|f\|\,\|x\|.
\]

\subsection{Сопряжённое пространство}

Множество всех линейных непрерывных функционалов на $X$ обозначается
\[
    X^*=B(X,\mathbb R)
\]
и называется \emph{сопряжённым}, или двойственным, пространством к $X$.
С операциями поточечного сложения и умножения на число и с нормой
\[
    \|f\|=\sup_{\|x\|\leq1}|f(x)|
\]
пространство $X^*$ является банаховым, даже если исходное нормированное
пространство $X$ не полно.

Двойственное пространство играет ключевую роль в определении слабой сходимости:
\[
    x_n\rightharpoonup x
    \quad\Longleftrightarrow\quad
    f(x_n)\to f(x)\quad\forall f\in X^*.
\]

\subsection{Теорема Хана--Банаха}

Одна из центральных теорем функционального анализа утверждает возможность
продолжения линейного функционала с подпространства на всё пространство.

\textbf{Теорема Хана--Банаха.}
Пусть $X$ --- нормированное пространство, $L\subset X$ --- линейное
подпространство и $f$ --- линейный ограниченный функционал на $L$. Тогда
существует линейный непрерывный функционал $F\in X^*$ такой, что
\[
    F|_L=f,
    \qquad
    \|F\|_{X^*}=\|f\|_{L^*}.
\]
То есть функционал можно продолжить на всё пространство без увеличения нормы.

Важное следствие: для любого $x_0\neq0$ существует $f\in X^*$ такое, что
\[
    \|f\|=1,
    \qquad
    f(x_0)=\|x_0\|.
\]
Следовательно, непрерывных линейных функционалов достаточно, чтобы различать
точки нормированного пространства.

\subsection{Теорема Рисса в гильбертовом пространстве}

Если $H$ --- гильбертово пространство, то его двойственное пространство имеет
особенно простой вид.

\textbf{Теорема Рисса о представлении функционала.}
Для любого $f\in H^*$ существует единственный элемент $u\in H$ такой, что
\[
    f(v)=(v,u),
    \qquad \forall v\in H,
\]
и при этом
\[
    \|f\|_{H^*}=\|u\|_H.
\]
Обратно, каждый фиксированный $u\in H$ задаёт таким способом линейный
непрерывный функционал. Поэтому $H$ и $H^*$ изометрически изоморфны.

Для пространства $L^2(\Omega)$ это означает, что любой непрерывный линейный
функционал имеет вид
\[
    f(v)=\int_\Omega v(x)\,u(x)\,dx
\]
для единственной функции $u\in L^2(\Omega)$.


\subsection{Доказательная схема эквивалентности непрерывности и ограниченности}

Если существует $C$ такое, что $|f(x)|\le C\|x\|$, то
\[
|f(x)-f(y)|=|f(x-y)|\le C\|x-y\|,
\]
поэтому $f$ непрерывен. Обратно, если $f$ непрерывен в нуле, то существует
$\delta>0$ такое, что из $\|x\|<\delta$ следует $|f(x)|<1$. Для $x\ne0$
применим это к
\[
    y=\frac{\delta}{2\|x\|}x.
\]
Получаем
\[
    |f(x)|\le \frac{2}{\delta}\|x\|.
\]
Таким образом, для линейных отображений достаточно проверять непрерывность только
в одной точке --- в нуле.

\subsection{Следствия теоремы Хана--Банаха}

Теорема Хана--Банаха имеет несколько важных следствий.
\begin{enumerate}
    \item Для каждого $x\ne0$ существует $f\in X^*$ такое, что
    \[
        \|f\|=1,\qquad f(x)=\|x\|.
    \]
    Поэтому
    \[
        \|x\|=\sup_{\|f\|\le1}|f(x)|.
    \]
    \item Если $M\subset X$ --- замкнутое линейное подпространство и
    $x_0\notin M$, то существует непрерывный линейный функционал, который
    обращается в нуль на $M$, но не в нуль на $x_0$. Это одна из форм
    теорем отделения.
\end{enumerate}

\subsection{Рисс и геометрия гильбертова пространства}

По теореме Рисса отображение
\[
    R:H\to H^*,\qquad Ru=(\,\cdot\,,u),
\]
является изометрическим изоморфизмом (в комплексном случае учитывается выбранная
конвенция линейности скалярного произведения). Поэтому всякий функционал можно
представлять вектором, а задачи на функционалы --- геометрически.

Например, если $M\subset H$ --- замкнутое подпространство и $x\in H$, то условие
ортогональной проекции
\[
    x-P_Mx\perp M
\]
можно эквивалентно записать как зануление соответствующего линейного функционала
на $M$.

\subsection{Слабая сходимость}

Последовательность $x_n$ слабо сходится к $x$, если
\[
    f(x_n)\to f(x)\qquad \forall f\in X^*.
\]
Слабая сходимость обычно слабее сходимости по норме. Если $x_n\to x$ по норме,
то автоматически $x_n\rightharpoonup x$. Обратное неверно в бесконечномерных
пространствах; стандартный пример --- ортонормированная последовательность в
гильбертовом пространстве.

\subsection{Что важно сказать на экзамене}

Линейный функционал --- это линейный оператор со значениями в поле чисел.
Для него непрерывность эквивалентна ограниченности, а норма равна
\[
    \|f\|=\sup_{\|x\|\leq1}|f(x)|.
\]
Все такие функционалы образуют двойственное пространство $X^*$. Теорема
Хана--Банаха позволяет продолжать функционал с подпространства без увеличения
нормы, а теорема Рисса в гильбертовом пространстве даёт представление
$f(v)=(v,u)$.

\newpage
